\section{A fixed-reference audit for distribution repair}
\label{sec:method}
The audit fixes a reference and regional targets, measures mass, spread,
and support, and tests a diagnosis-guided intervention on independent
outcomes (Fig.~\ref{fig:pipeline}; Supp.~Algorithm~1).

\subsection{Setting and notation}
\label{sec:setting}
Let $\mathcal X$ be the image space. A configuration $a$ specifies the
checkpoint, schedule, guidance, solver, and randomization of a generator,
and $Q_a$ denotes its output law on $\mathcal X$. The audit compares $Q_a$
with a reference law $P$. We use the law $P_T$ of a resolved teacher to
measure sampler-relative change, or the observed-data law $P_D$ to measure
data-relative representation. Before evaluation we declare the
representation, regions, effect sizes, fidelity margins, comparison
family, and budget.

Disjoint reference pools fit the regions and support score
($R_{\rm fit}$), calibrate the threshold ($R_{\rm cal}$), and estimate
reference outcomes ($R_{\rm eval}$). Candidate images enter neither
fitting nor calibration. A diagnosis pool selects the intervention;
an independent outcome pool evaluates the frozen comparison. The audit
reports regional decisions and intervals, conditional spread, fidelity,
and cost.

\subsection{Representation, regions, and pooled spread}
\label{sec:regions}
A representation $\phi:\mathcal X\to\mathbb R^d$ maps an image to $d$
features. Pixel, latent, DINOv2, and CLIP features offer different views
of the images~\cite{oquab2024dinov2,radford2021clip}. Principal component
analysis (PCA) fitted on $R_{\rm fit}$ gives a mean $m\in\mathbb R^d$,
an orthonormal basis $U\in\mathbb R^{d\times k}$ of $k$ retained
directions, and coordinates $z(x)=U^\top(\phi(x)-m)\in\mathbb R^k$.
For feature regions, we scale each coordinate of $z$ by its reference-fit
standard deviation and fit $K$ $K$-means cells to the scaled coordinates.
The regions use these scaled coordinates, the spread ratio of
\cref{eq:spread} uses $z$ itself, and the support score of
\cref{eq:support} uses whitened coordinates. The fixed assignment
$J(x)\in\{1,\ldots,K\}$ defines the regions $A_j=\{x:J(x)=j\}$, which
partition $\mathcal X$.

For semantic regions, the same classifier and rejection rule are applied
to both laws. The partition includes rejection, so rejected images stay
in the denominator of each region's probability. Held-out real accuracy
validates the labeling rule on real images. Feature fidelity and
uncurated samples help assess its interpretation on generated images.

Write the regional masses as $p_j=P(X\in A_j)$ and
$q_{aj}=Q_a(X\in A_j)$. Under $P$, let $Y=\phi(X)$ have mean $\mu$,
with conditional mean $\mu_j$ and covariance $\Sigma_j$ in $A_j$.
The law of total covariance gives
\begin{equation}
\operatorname{Cov}_{P}(Y)
=\sum_j p_j\Sigma_j+\sum_j p_j(\mu_j-\mu)(\mu_j-\mu)^\top,
\label{eq:totalcov}
\end{equation}
with zero-mass terms contributing zero. Pooled spread combines
within-region variation, regional separation, and their probability
weights; the audit measures mass and conditional spread separately.

\subsection{Certifying a regional change in mass}
\label{sec:mass}
To certify a regional change, we fix all events before drawing evaluation
images independently within each law. Reference and candidate draws may
be paired across laws. The guarantee assumes fixed evaluation sizes.
Pooling samples drawn under different prompts requires a sampling model
for the prompt mixture, so we report prompt-level uncertainty separately.

Let $M$ count all declared region, condition, configuration, and
representation comparisons, and let $\alpha_{\rm occ}$ be the family
error budget. Let $[p_j^-,p_j^+]$ and $[q_{aj}^-,q_{aj}^+]$ be exact
binomial intervals~\cite{clopper1934confidence} with per-tail error
$\varepsilon=\alpha_{\rm occ}/(4M)$.

\paragraph{Proposition 1 (simultaneous ratio intervals).}
\emph{With probability at least $1-\alpha_{\rm occ}$, every positive-mass
region in the family satisfies}
\begin{equation}
q_{aj}/p_j\in\left[q_{aj}^-/p_j^+,\ q_{aj}^+/p_j^-\right].
\label{eq:occupancy}
\end{equation}

\noindent\emph{Proof.} A union bound over the $4M$ one-sided endpoints
gives simultaneous coverage at least $1-\alpha_{\rm occ}$; monotonicity of the
ratio gives \cref{eq:occupancy}. \hfill$\square$

For $p_j>0$, the upper endpoint is infinite when $p_j^-=0$. Only eligible
regions receive a ratio decision. An eligible region has
$p_j^-\geq\pi_{\min}$, a minimum reference mass. An upper endpoint below
the ratio threshold $\rho$ certifies \emph{depletion}; a lower endpoint
above $\rho$ certifies \emph{retention} at that resolution. Otherwise the
test is unresolved. For signed mass balance, we bound $q_{aj}-p_j$ by
$[q_{aj}^--p_j^+,q_{aj}^+-p_j^-]$ with a margin $\delta>0$. The interval
must lie below, above, or within $[-\delta,\delta]$ to certify depletion,
excess, or equivalence, respectively.

\paragraph{Corollary 1 (samples needed).}
\emph{Suppose a region is eligible, $0<\rho p_j^-<1$, and none of $n_Q$
candidate draws falls in it. Then \cref{eq:occupancy} certifies depletion
if and only if}
\begin{equation}
n_Q>\frac{\log\varepsilon}{\log(1-\rho p_j^-)}.
\label{eq:planning}
\end{equation}

\noindent\emph{Proof.} With zero candidate hits,
$q_{aj}^+=1-\varepsilon^{1/n_Q}$; solving
$q_{aj}^+<\rho p_j^-$ gives \cref{eq:planning}. \hfill$\square$

The bound makes the sample cost of rare regions and broad comparison
families explicit.

\subsection{Conditional spread}
\label{sec:spread}
To measure variation separately from mass, we compare the covariance
traces of the unwhitened coordinates within $A_j$:
\begin{equation}
r_{aj}=\frac{\operatorname{tr}\operatorname{Cov}_{Q_a}[z(X)\mid A_j]}
{\operatorname{tr}\operatorname{Cov}_{P}[z(X)\mid A_j]}.
\label{eq:spread}
\end{equation}
This ratio is defined when the region has positive mass under both laws
and the reference trace is positive. The trace sums coordinate variances
after centering each law separately. Approximate influence-function
intervals for $\log r_{aj}$ lying wholly below or above declared margins
flag contraction or expansion. For text-to-image generation, we divide
the mean within-prompt trace of the candidate by that of the reference
in each representation and bootstrap whole matched prompts.

\subsection{Reference support}
\label{sec:support}
A support score $S(x)$ combines the whitened distance to the nearest
reference-fit point and the scaled PCA residual norm in quadrature:
\begin{equation}
S(x)=\sqrt{d_{\rm NN}(x)^2+
\bigl(\|\phi(x)-m-Uz(x)\|_2/s_{\rm res}\bigr)^2},
\label{eq:support}
\end{equation}
where $d_{\rm NN}(x)$ is the nearest fit-sample distance in whitened PCA
coordinates and $s_{\rm res}$ is the residual scale (Supp.~Sec.~C.1).
The residual captures departures in directions discarded by projection.
With $n_{\rm cal}$ exchangeable reference calibration samples, the
threshold $\tau$ is the $\lceil(n_{\rm cal}+1)0.95\rceil$-th score in
ascending order~\cite{vovk2013conditional}. If this rank is unavailable,
the gate is unresolved. Independent evaluations estimate the excess tail
probability $Q_a(S>\tau)-P(S>\tau)$, which exact binomial endpoints bound
(Supp.~Sec.~A.3). The gate passes if the upper bound is at most a declared
tolerance and fails if the lower bound exceeds it. Otherwise it is
unresolved. Expansion with a failed support gate is off-reference expansion.

\subsection{From diagnosis to a controlled intervention}
\label{sec:intervene}
For the same semantic event, let $p_{D,j}=P_D(A_j)$ and $p_{T,j}=P_T(A_j)$.
When both are positive,
\begin{equation}
\frac{q_{aj}}{p_{D,j}}=\frac{q_{aj}}{p_{T,j}}\cdot\frac{p_{T,j}}{p_{D,j}}
\label{eq:decomp}
\end{equation}
separates sampler--teacher from teacher--data mass differences.
Teacher-relative deficits motivate a starting-noise change at fixed
denoiser, grid, and call budget, or a solver/grid comparison under the
same conditioning. A validated data-relative imbalance persisting across
qualified samplers motivates ordinary versus targeted continuation from
the same parent at equal update budgets. After diagnosis, we freeze one
change and its control, then evaluate the targeted outcome, fidelity,
coverage, and cost on independent samples.

\paragraph{Why the starting covariance matters.}
Common schedules with leading timestep spacing start sampling from a
noise level that the standard Gaussian does not match~\cite{lin2024flawed},
and data-dependent priors change the starting distribution through
training~\cite{lee2022priorgrad}. We quantify the effect of the mismatch
and test a training-free correction. Consider a centered clean Gaussian
coordinate $X_0$ with variance
$\lambda>0$. Variance-preserving noising gives
$X_t=\sqrt{\bar\alpha_t}X_0+\sqrt{1-\bar\alpha_t}\,\epsilon$ at
timestep $t$, with cumulative signal coefficient $\bar\alpha_t\in[0,1]$
and independent $\epsilon\sim\mathcal N(0,1)$. Write
$v_t=\bar\alpha_t\lambda+1-\bar\alpha_t$ for the noisy variance and
$u_t=(\sqrt{\bar\alpha_t\lambda},\sqrt{1-\bar\alpha_t})$.
Let $X_{\rm out}$ be the final coordinate, and let $\nu_{\rm start}$ and
$v_{\rm start}$ be the actual and matched starting variances.

\paragraph{Proposition 2 (initialization and integration).}
\emph{With the conditional mean noise predictor
$\mathbb E[\epsilon\mid X_t]$ and no clipping, a deterministic DDIM
transition to an earlier timestep $s$ multiplies the coordinate by
$\sqrt{v_s/v_t}\cos\angle(u_s,u_t)$. On a fixed grid, the variance factors
telescope to the clean endpoint:}
\begin{equation}
\frac{\operatorname{Var}(X_{\rm out})}{\lambda}
=\underbrace{\frac{\nu_{\rm start}}{v_{\rm start}}}_{\text{initialization}}
\underbrace{\prod_{t\rightarrow s}\cos^2\angle(u_s,u_t)}_{\text{integration}},
\label{eq:gaussian}
\end{equation}
\emph{where the product runs over adjacent grid transitions.}

The derivation builds on exact Gaussian analyses
\cite{benita2025spectral,pierret2025gaussian}
(Supp.~Secs.~D--F). Matching the noisy marginal sets the initialization
factor to one; refining a grid with fixed endpoints improves the
integration factor. Changing the leading-grid budget also changes the
starting timestep. \Cref{eq:gaussian} extends directionwise to shared
eigenvectors.

For a leading grid starting at timestep $t_0$, the tested initializer is
\begin{equation}
X_{\rm start}=
\big[(1-\bar\alpha_{t_0})I+\bar\alpha_{t_0}\widehat C\big]^{1/2}Z,
\qquad Z\sim\mathcal N(0,I),
\label{eq:covariance_start}
\end{equation}
where $X_{\rm start}$ is the initial noisy image, $Z$ is standard Gaussian
noise, $I$ is the identity in pixel space, $\widehat C$ is an estimate of
the clean covariance, and the square root is symmetric. We call this the
\emph{covariance start} and the usual $X_{\rm start}=Z$ the
\emph{standard start}. The covariance start adds no denoiser calls.

\subsection{Implementation details}
\label{sec:implementation}
We retain up to $k=32$ PCA directions and fit $K=10$ $K$-means cells.
Depletion uses $\rho=0.5$ and $\pi_{\min}=0.01$ by default, and the digit
study uses $\delta=0.03$ and charges the shared reference bounds once
(Supp.~Sec.~C.8). A finite estimated log-ratio spread interval requires at
least 30 samples from each law and positive sample trace in both;
intervals wholly below 0.9 or above 1.1 flag contraction or expansion.
The residual scale $s_{\rm res}$ is the median fit residual, floored at
$10^{-6}$, and the support gate tolerance is 0.05. The feature-only
default gives occupancy and support error budgets of 0.025 each. The
approximate spread intervals are outside that exact error guarantee.
Signed-balance and routing follow-ups allocate 0.025/0.0125/0.0125 to
mass/fidelity/support. CIFAR selection also checks teacher precision,
text-to-image checks alignment and learned preference, and digits use
classifier consensus and feature precision. For CIFAR-10, $I$ is the
identity in 3,072-dimensional pixel space, $\widehat C$ uses 128
teacher-fit pixel components and an isotropic residual, and the two-call
grid starts at $t_0=500$.
